%%%%%%%%%%%%%%%%%%%%%%%%%%%%%%%%%
%
%       EXERCÍCIO
%
%%%%%%%%%%%%%%%%%%%%%%%%%%%%%%%%%

\ifdefstring{\atividade}{prova}{%
    \ifdefstring{\modo}{objetivo}{%
        \renewcommand{\valorquestao}{\ValorQObj\ ponto}
    }{%
        \renewcommand{\valorquestao}{\ValorQDisc\ pontos}
    }%
}%

\begin{exercicioBanco}[\valorquestao]
Seja \(f\) uma função afim tal que \(f(1) = 2\) e \(f(3) = 5\). Determine a lei que representa essa função.

% Define as alternativas
\newcommand{\alternativas}{%
    \begin{center}
        \begin{tabularx}{\textwidth}{XXX}
            (a) \(f(x) = 2x - 1\). &
            (b) \(f(x) = \dfrac{3}{2}x - \dfrac{1}{2}\). &
            (c) \(f(x) = x + 1\). \\[10pt]
            (d) \(f(x) = \dfrac{3}{2}x + 1\). &
            (e) \(f(x) = \dfrac{3}{2}x + \dfrac{1}{2}\).
        \end{tabularx}
    \end{center}
}

% Resposta correta
\newcommand{\resposta}{E}

% Lógica condicional para exibição
\ifdefstring{\atividade}{lista}{%
    \alternativas
    \vspace{0.5em}
    
    \noindent\textbf{Resposta:} letra \textbf{\resposta}.
}{%
    \ifdefstring{\modo}{objetivo}{%
        \alternativas
    }{%
        % modo = discursiva → não mostra alternativas
    }
}
\end{exercicioBanco}

